%% kapa: exam=1 round=20 year=2009 period=1 form=A subject=civil_law subject_ko=민법 seq=13 number=13
%% source: Q-Net 감정평가사 기출문제 (cst00309) artlSeq=5024178 file=extracted/1008/5024178/2009년 제20회 감정평가사 1차시험 기출문제/1교시.hwp
%% conversion: kordoc 4.24.51 -> tools/gen1.py (auto); review=proofread-sonnet; flags=none
\begin{question}{13}
권리능력 없는 사단(법인 아닌 사단)에 관한 설명으로 옳은 것은? (다툼이 있으면 판례에 의함)
\begin{choices}
  \item 사단법인에 관한 민법의 규정 중 법인격을 전제로 하는 것을 제외하고는, 이를 권리능력 없는 사단에 유추적용할 수 있다.
  \item 권리능력 없는 사단의 재산소유형태는 합유이다.
  \item 종중은 공동선조와 성(姓)과 본(本)을 같이하는 성년 남자로만 구성된다.
  \item 권리능력 없는 사단은 소송상 당사자능력이 없다.
  \item 권리능력 없는 사단인 교회의 교인들의 일부가 집단적으로 탈퇴하거나 소속교단을 변경하면 교회가 분열된 것으로 본다.
\end{choices}
\end{question}
\answer{1}
